\section{Conclusiones y Trabajo Futuro}\label{sec:conclusions}

\subsection{Conclusiones técnicas}
Este trabajo desarrolló y evaluó un pipeline reproducible de minería de datos
para la predicción binaria de heladas en el territorio peruano, sobre un conjunto
de $1\,672\,827$ observaciones meteorológicas. De la evidencia experimental se
extraen las siguientes conclusiones directas:

\begin{enumerate}
\item El ensamble Random Forest optimizado alcanzó un \emph{Recall} de $92.34\%$
en la detección de heladas y un ROC-AUC de $0.9690$, con un índice Kappa de Cohen
de $0.6710$ (concordancia sustancial), validando el sistema como una herramienta
de alerta temprana de alta sensibilidad.
\item La estrategia de ponderación balanceada de clases combinada con
optimización directa sobre F1-Score Macro resultó eficaz para manejar un
desbalance severo ($16.24\%$ frente a $83.76\%$) sin recurrir a generación
sintética de muestras, evitando así la amplificación de la complejidad de datos
ya afectados por ruido de instrumentación.
\item La cercanía entre el F1-Score Macro de validación interna ($0.8330$) y el de
prueba ($0.8338$) evidencia una capacidad de generalización estable y ausencia de
sobreajuste significativo dentro del dominio del conjunto evaluado.
\item El pipeline de escalamiento robusto (\texttt{RobustScaler}) seguido de PCA
redujo la dimensionalidad de 15 a 5 componentes reteniendo el $95\%$ de la
varianza, demostrando que el problema admite una representación compacta sin
pérdida sustancial de capacidad predictiva.
\end{enumerate}

\subsection{Lecciones aprendidas sobre los datos y la ingeniería de características}
El proceso de desarrollo dejó tres aprendizajes de valor sobre el comportamiento
de los datos y su preparación.

Primero, el diagnóstico visual exploratorio previo a cualquier transformación
matemática resultó indispensable: el código de error de sensor \texttt{-999} se
encontraba enmascarado y no habría sido detectado por una inspección estructural
convencional (\texttt{df.info()} reportaba cero nulos). Su omisión habría
distorsionado irreversiblemente tanto la matriz de correlación como el
escalamiento, comprometiendo todo el análisis posterior. Esto confirma que un
conjunto declarado ``preprocesado'' por su fuente exige, no obstante, una
auditoría de calidad independiente.

Segundo, la ingeniería de la variable objetivo demostró que la prevención de fuga
de información es un requisito de diseño y no un detalle: conservar
\texttt{temp\_min\_2m} entre los predictores habría inducido una regla
determinista trivial con una correlación artificial cercana a $-1$, ocultando las
verdaderas relaciones climáticas del sistema.

Tercero, el análisis de impureza de Gini reveló que la capacidad discriminante no
reside en la dirección de mayor varianza global (PC1), sino en componentes
latentes de varianza intermedia (PC3, PC2). Esta lección advierte contra la
suposición ingenua de que las primeras componentes principales concentran
siempre la señal predictiva más relevante.

\subsection{Trabajo futuro}
Se identifican cuatro líneas de trabajo futuro de carácter metodológico.

\emph{Validación estadística rigurosa.} La línea prioritaria consiste en dotar al
estudio de respaldo inferencial mediante la incorporación de modelos de línea base
(Regresión Logística, SVM, XGBoost) entrenados bajo el mismo pipeline, ejecutando
validación cruzada estratificada repetida con múltiples semillas y aplicando
pruebas de significancia (Friedman con post-hoc de Nemenyi, o Wilcoxon con
corrección de Bonferroni) para contrastar formalmente el modelo propuesto frente a
las alternativas.

\emph{Enriquecimiento de características espaciotemporales.} Se propone incorporar
descriptores que capturen la dinámica temporal del fenómeno, en particular una
codificación trigonométrica (seno/coseno) de las variables cíclicas de día y mes,
así como reformular el problema como un pronóstico a $t+k$ horas en lugar de una
clasificación sobre observaciones concurrentes, acercando el sistema a un
escenario operativo real de alerta anticipada.

\emph{Exploración de arquitecturas profundas.} Dada la escala del conjunto (más de
1.6 millones de observaciones), resulta pertinente evaluar arquitecturas de
aprendizaje profundo ---redes recurrentes (LSTM/GRU) o convolucionales
temporales--- capaces de modelar dependencias secuenciales, contrastándolas con el
ensamble Random Forest como línea base robusta.

\emph{Técnicas de rebalanceo y aprendizaje semisupervisado.} Se plantea explorar
técnicas de generación sintética adaptativa (ADASYN) sobre el espacio de
entrenamiento para intentar reducir la tasa de falsos positivos sin sacrificar la
sensibilidad alcanzada, así como estrategias de aprendizaje semisupervisado que
aprovechen registros meteorológicos no etiquetados de estaciones adicionales para
mejorar la generalización geográfica del modelo.